\documentclass[11pt]{article}

\usepackage[utf8]{inputenc}
\usepackage[T1]{fontenc}
\usepackage[spanish]{babel}
\usepackage[a4paper,margin=2cm]{geometry}
\usepackage{amsmath}
\usepackage{graphicx}
\usepackage{pgfplots}
\usepackage{url}

\graphicspath{{./}{docs/}}

\begin{document}

\begin{titlepage}
    \centering
    \vspace*{2.5cm}
    \includegraphics[width=0.44\textwidth]{logo_tec.png}\par
    \vspace{1.5cm}
    {\Large\bfseries Informe técnico\par}
    \vspace{0.65cm}
    {\large Procesamiento paralelo de imágenes\\con Scheme y Erlang\par}
    \vspace{0.8cm}
    {\large Lenguajes de Programación\par}
    \vspace{1.1cm}
    {\large Estudiante:\par}
    \vspace{0.15cm}
    {\large Daniel Araya Pérez\\Carné: 2025078995\par}
    \vspace{1.1cm}
    {\large Profesor:\par}
    \vspace{0.15cm}
    {\large Eddy Ramírez Jiménez\par}
    \vfill
    {\large Setiembre de 2026\par}
    \vspace*{1.5cm}
\end{titlepage}

\section{Resumen}

El proyecto combina Erlang para coordinar procesos y Scheme para transformar
imágenes PPM P3. Erlang divide la imagen, ejecuta una
instancia de Scheme por región y reúne los resultados. Se implementaron
difuminado gaussiano y escala de grises. El rendimiento se midió con 1, 2, 4
y 8 procesos.

\section{Diseño de la solución}

\subsection{Arquitectura}

El archivo \texttt{image\_processor} inicia el programa. En Erlang,
\texttt{main.erl} recibe los parámetros, \texttt{ppm.erl} lee la imagen,
\texttt{partition.erl} crea las regiones y \texttt{coordinator.erl} dirige
la ejecución. Cada proceso de \texttt{worker.erl} usa
\texttt{scheme\_bridge.erl} para ejecutar su propia instancia de
\texttt{scheme/worker.scm}. Al terminar, el coordinador ordena las regiones
y escribe la imagen completa.

Una imagen se representa mediante su ancho, su alto y una lista de filas.
Cada píxel tiene tres componentes: rojo, verde y azul. Los lectores y
escritores de Erlang y Scheme usan PPM P3, la variante de texto del formato
\cite{netpbm-ppm}, con valor máximo 255. Los módulos Scheme usan Full Swindle
\cite{racket-swindle}.

\subsection{División de la imagen}

La imagen se divide en franjas horizontales de ancho completo. Para una altura
$H$ y $p$ procesos, cada franja recibe al menos $\lfloor H/p \rfloor$ filas.
Las $H \mathbin{\mathrm{mod}} p$ filas restantes se reparten, una por una, entre las primeras
franjas. Por eso, la cantidad de procesos no puede superar la altura.

El difuminado necesita filas vecinas de otras franjas. Para un kernel de tamaño
impar $k$, cada región recibe hasta $r=(k-1)/2$ filas de contexto por arriba y
por abajo. Estas filas forman el halo. El proceso devuelve solamente sus filas
asignadas. La escala de grises no necesita halo. Como las franjas ocupan todo
el ancho, no se agregan columnas de contexto.

\subsection{Concurrencia y reconstrucción}

El coordinador crea un proceso Erlang por región con \texttt{spawn\_monitor}
\cite{erlang-processes}.
Cada proceso envía un mensaje \texttt{region\_completed} al terminar. La
recepción ocurre conforme llegan los mensajes, sin exigir que los procesos
terminen en el orden de las franjas. Después, Erlang ordena las regiones por
su fila inicial y concatena sus filas para reconstruir la imagen.

\section{Comunicación entre Erlang y Scheme}

El intercambio usa archivos temporales y argumentos de línea de comandos.
Para cada región, Erlang crea una carpeta única en \texttt{/tmp} con tres
rutas: entrada, salida y estado. Escribe la región, incluido su halo, como
PPM P3 y ejecuta \texttt{scheme/worker.scm} mediante el comando
\texttt{racket}.

Se eligieron archivos para reutilizar los lectores y escritores PPM de ambos
lenguajes sin gestionar conexiones. Las carpetas separadas permiten ejecutar
trabajadores en paralelo sin compartir rutas, a cambio de más operaciones
de disco.

Los argumentos indican las tres rutas, el filtro y los límites de las filas
que deben devolverse, expresados respecto de la región con halo. Para
\texttt{gaussian} también incluyen el tamaño del kernel y sigma. Scheme lee
la entrada, aplica el filtro, escribe solo la parte asignada y al final escribe
\texttt{ok} en el archivo de estado.

Erlang acepta el resultado solo si encuentra ese estado. Después lee la
salida, comprueba sus dimensiones y elimina los archivos temporales del
trabajador. Cada proceso Erlang espera a su propia instancia de Scheme, pero
los distintos procesos Erlang trabajan de forma concurrente.

\newpage

\section{Filtros y bordes}

\subsection{Difuminado gaussiano}

\texttt{gaussian\_blur.scm} construye el kernel a partir de su tamaño y sigma.
El peso de cada posición $(x,y)$ es proporcional a

\begin{equation}
    w(x,y)=\exp\left(-\frac{x^2+y^2}{2\sigma^2}\right).
\end{equation}

Los pesos se dividen entre su suma para que el total sea uno
\cite{mit-gaussian}.
\texttt{convolution.scm} combina cada componente RGB con esos pesos, redondea
el resultado y lo limita al intervalo de 0 a 255. El tamaño predeterminado
es $3\times3$ y el sigma predeterminado es $0.8493$. También se aceptan
tamaños impares desde 3 y valores positivos de sigma.

Cuando una posición del kernel queda fuera de la imagen, se usa el píxel
más cercano del borde. Scheme aplica esta regla tanto en filas como en
columnas. El halo proporciona los píxeles reales de las franjas vecinas, de
modo que un límite entre procesos no se trata como borde de la imagen.

\subsection{Escala de grises}

\texttt{grayscale.scm} calcula

\begin{equation}
    Y=\operatorname{round}(0.299R+0.587G+0.114B)
\end{equation}

y produce el píxel $(Y,Y,Y)$. MathWorks presenta estos coeficientes en su
documentación sobre conversión de RGB a gris \cite{mathworks-grayscale}.
Este filtro no usa kernel ni sigma.

\section{Manejo de errores}

Erlang rechaza parámetros inválidos, como una cantidad de procesos no
positiva, un kernel par o rutas de entrada y salida iguales. Si Scheme falla
antes de escribir \texttt{ok}, Erlang detecta que falta el estado de éxito.
El trabajador termina con error y el monitor informa al coordinador.

Ante el fallo de un trabajador, el coordinador cancela los pendientes e
informa el problema. La imagen final solo se escribe cuando llegaron todas
las regiones válidas, por lo que un fallo no genera una imagen incompleta.
No se implementó la repetición automática de una región fallida.

\section{Evaluación del rendimiento}

\subsection{Método}

Se usó \texttt{samples/sample\_1.ppm}, de $640\times413$ píxeles, con el
difuminado predeterminado: kernel $3\times3$ y sigma $0.8493$. Las mediciones
se hicieron en Ubuntu bajo WSL, con 12 procesadores lógicos, Racket 9.3 y el
emulador Erlang 13.2.2.5. Se realizaron ocho corridas con 1 y 2 procesos y
dieciséis con 4 y 8 procesos. La tendencia general se repitió entre corridas.
El tiempo informado por el programa incluye lectura, división, procesamiento,
reconstrucción y escritura de la imagen. No incluye la compilación previa
del lanzador ni el arranque de la máquina virtual Erlang.

Se toma la mediana de las mediciones como $T_p$ para reducir el efecto de
valores aislados. La aceleración y la eficiencia se calculan mediante

\begin{equation}
    S_p=\frac{T_1}{T_p},
    \qquad
    E_p=\frac{S_p}{p}.
\end{equation}

\begin{table}[htbp]
    \centering
    \begin{tabular}{r|r|r|r|r}
        \hline
        Procesos & Corridas & $T_p$ (s) & $S_p$ & $E_p$ (\%) \\
        \hline
        1 & 8 & 4.228 & 1.00 & 100.0 \\
        2 & 8 & 2.234 & 1.89 & 94.6 \\
        4 & 16 & 1.682 & 2.51 & 62.8 \\
        8 & 16 & 1.504 & 2.81 & 35.1 \\
        \hline
    \end{tabular}
    \caption{Medianas de tiempo, aceleración y eficiencia}
    \label{tab:rendimiento}
\end{table}

\begin{figure}[htbp]
    \centering
    \begin{tikzpicture}
        \begin{axis}[
            ybar,
            ymin=0,
            ymax=5,
            width=0.8\textwidth,
            height=6cm,
            xlabel={Procesos},
            ylabel={Tiempo (s)},
            symbolic x coords={1,2,4,8},
            xtick=data,
            nodes near coords
        ]
            \addplot coordinates {(1,4.228) (2,2.234) (4,1.682) (8,1.504)};
        \end{axis}
    \end{tikzpicture}
    \caption{Tiempo mediano en segundos por cantidad de procesos}
    \label{fig:tiempo}
\end{figure}

\begin{figure}[htbp]
    \centering
    \begin{tikzpicture}
        \begin{axis}[
            ybar,
            ymin=0,
            ymax=3.5,
            width=0.8\textwidth,
            height=6cm,
            xlabel={Procesos},
            ylabel={Aceleración},
            symbolic x coords={1,2,4,8},
            xtick=data,
            nodes near coords
        ]
            \addplot coordinates {(1,1.00) (2,1.89) (4,2.51) (8,2.81)};
        \end{axis}
    \end{tikzpicture}
    \caption{Aceleración por cantidad de procesos}
    \label{fig:aceleracion}
\end{figure}

\subsection{Análisis}

La mediana baja de 4.228 segundos con un proceso a 1.504 segundos con ocho.
La mayor aceleración medida es $2.81$, con ocho procesos. La mejora de
cuatro a ocho procesos es menor que la de uno a dos. Algunas corridas
fueron más lentas, pero la tendencia general se mantuvo.

El tiempo no se dedica solo al filtro. Cada región requiere crear un proceso
Scheme, escribir y leer archivos PPM, y copiar datos entre Erlang y Scheme.
También se pagan la división y la reconstrucción. Estos costos reducen la
eficiencia al aumentar $p$. No se midió cada etapa por separado.

\section{Limitaciones}

\begin{itemize}
    \item Solo se aceptan imágenes PPM P3 válidas, con valor máximo 255 y sin
    comentarios.
    \item Cada región inicia una instancia nueva de Scheme y usa archivos
    temporales, lo que agrega trabajo a cada ejecución.
    \item Un trabajador fallido se informa, pero su región no se vuelve a
    procesar automáticamente.
    \item Si se cancela un trabajador mientras usa su carpeta temporal, pueden
    quedar archivos en \texttt{/tmp}.
    \item Las rutas se colocan entre comillas simples en un comando de shell.
    Una ruta que contenga una comilla simple puede fallar.
\end{itemize}

\section{Conclusiones}

Erlang resulta adecuado para coordinar las regiones porque permite crear
procesos concurrentes, recibir sus resultados por mensajes y detectar fallos
con monitores. Scheme resulta adecuado para los filtros porque sus funciones
recursivas procesan listas de filas y píxeles, y construyen el kernel sin
modificar la imagen de entrada.

Los halos conservan los vecinos necesarios en los límites entre regiones.
En la imagen medida, ocho procesos dieron el menor tiempo, pero la eficiencia
bajó al aumentar la cantidad de procesos por el trabajo adicional de
coordinación y comunicación.

\clearpage

\begin{thebibliography}{9}
    \bibitem{erlang-processes}
    Erlang, \emph{Processes}, Erlang System Documentation.
    \url{https://www.erlang.org/doc/system/ref_man_processes.html}.

    \bibitem{racket-swindle}
    Racket, \emph{Swindle}, Racket Documentation.
    \url{https://docs.racket-lang.org/swindle/index.html}.

    \bibitem{netpbm-ppm}
    Netpbm, \emph{PPM Format Specification}.
    \url{https://netpbm.sourceforge.net/doc/ppm.html}.

    \bibitem{mit-gaussian}
    A. Torralba, P. Isola y W. Freeman,
    \emph{Foundations of Computer Vision}, sección 17.3, Gaussian Filter.
    \url{https://visionbook.mit.edu/blurring_2.html}.

    \bibitem{mathworks-grayscale}
    MathWorks, \emph{rgb2gray: Convert RGB Image or Colormap to Grayscale},
    sección Algorithms.
    \url{https://www.mathworks.com/help/matlab/ref/rgb2gray.html}.
\end{thebibliography}

\end{document}
